\section{Informative instruments without posterior contraction}\label{sec:effective-realizations}
The preceding classical and quantum realizations yield explicit geometric rates because the current task response is report-recoverable. The exact-budget construction has a different, strictly broader hypothesis. We now verify it on informative experiments in which the current report does not determine the posterior. Neither proof below invokes a realization-specific theorem to establish the general construction.

\subsection{Controlled Gaussian hidden-state experiments}
Let the hidden state set be $\{1,\ldots,d\}$ and let the action set be finite. An action first makes a transition $P_a^\theta(j\mid i)$ and then reports a vector $y\in\R^p$ with density $g_{aj}^\theta(y)$ of $\mathcal N(m_{aj}^\theta,\Sigma_{aj}^\theta)$. Thus the raw joint kernel and its unnormalized instrument are
\begin{equation}\label{eq:effective-gaussian}
 K_a^\theta(j,dy\mid i)=P_a^\theta(j\mid i)g_{aj}^\theta(y)\,dy,
 \qquad L_a^\theta(y)_{ji}
       =P_a^\theta(j\mid i)g_{aj}^\theta(y)/w(y),
\end{equation}
where $\lambda(dy)=w(y)dy$ is a fixed nondegenerate Gaussian probability with explicitly known density, chosen independently of $\theta$. The ordered physical space is $\R^d$ with the nonnegative cone and sum as mass. A prepared hidden state has distribution $x_0^\theta$. After stopping, an independent protected audit returns the one-hot vector of the current hidden state. Its mean is the posterior vector and its second moment is one. Decoders lie in the simplex, and the squared loss is bounded by $2$.

\begin{theorem}[Exact-budget finite construction for informative Gaussian kernels]\label{thm:effective-gaussian}
Suppose the family of $x_0^\theta,P_a^\theta,m_{aj}^\theta,\Sigma_{aj}^\theta$ is effectively compact and computable, with
\[
 \norm{m_{aj}^\theta}\le C_m,\qquad
 c_\Sigma I\le\Sigma_{aj}^\theta\le C_\Sigma I,
 \qquad 0<c_\Sigma\le C_\Sigma<\infty.
\]
Assume the supplied report digitization can implement bounded-box partitions with a certified vanishing boundary error, and records the entire complement as a charged tail outcome. Then Theorems~\ref{thm:main}(i) and \ref{thm:effective-budget} hold for every rational $N\ge1$ and $M\ge2$.
No positivity floor, invertibility, irreducibility, or mixing rate is required for $P_a^\theta$. The theorem covers identity and reducible transitions and unknown fixed parameters using one common program.
\end{theorem}
\begin{proof}
Each matrix in \eqref{eq:effective-gaussian} is positive, its integrated column sums are one, and its integral operator norm is finite. The protected one-hot audit gives every coordinate of the posterior as an executable future-test evaluation. These evaluations separate the simplex. All unnormalized report pullbacks remain in the $d$-dimensional linear space. This verifies the finite predictive representation from the raw experiment. If a smaller test family is used, its invariant separating span, rather than the whole simplex, is the relevant quotient. Here the audit already separates all coordinates.

We verify the uniform approximation certificate. For $R\ge2C_m+1$, Gaussian tails and the covariance bounds give computable constants $C,c>0$ such that
\begin{equation}\label{eq:gaussian-tail-compile}
 \sup_{\theta,a,j}\int_{\norm y_\infty>R}g_{aj}^\theta(y)dy
        \le C\exp(-cR^2).
\end{equation}
For example, apply the one-dimensional Gaussian Chernoff bound to each coordinate and take the union bound. Partition the box $[-R,R]^p$ into half-open rational cubes of side at most $h$, and keep its complement as one additional cell. On the box, $w$ has a positive computable minimum. Differentiating the Gaussian formula shows that $g_{aj}^\theta/w$ has a computable uniform gradient bound $C_R$. The mean-value theorem bounds the oscillation of every instrument entry on a cube by $C_R\sqrt p\,h$. In a fixed finite dimension, summing entry bounds bounds the induced $\ell^1$ operator norm, so the integrated interior approximation error is at most $C'_Rh$.

On the tail cell, positivity gives
\[
 \int_{\mathrm{tail}}\bnorm{L-\overline L}\,d\lambda
    \le2\int_{\mathrm{tail}}\bnorm L\,d\lambda
    \le C'\exp(-cR^2).
\]
For the last inequality, the maximum column sum is at most the sum of all entries, bounded by a fixed-dimensional sum of Gaussian tails. Consequently
\begin{equation}\label{eq:gaussian-partition-modulus}
 \eta_{R,h}\le C'_Rh+C'\exp(-cR^2).
\end{equation}
First choose $R$ and then $h$. The original tail remains part of every physical call and of every cell integral; it is neither discarded nor conditioned away. A report-digitization error is added to this bound as specified in Definition~\ref{def:effective}.

For model continuity, derivatives of a Gaussian density with respect to its mean and covariance are that density times a polynomial of degree at most two in $y-m$, with coefficients bounded by powers of $c_\Sigma^{-1}$. Their absolute integrals are uniformly bounded on the stated parameter body. Integrate these derivatives along a straight segment in mean--covariance space; the covariance segment retains the eigenvalue bounds. This proves a computable $L^1$ modulus for the densities. Matrix multiplication by stochastic $P_a^\theta$ and finite-dimensional triangle inequalities give the required integral instrument modulus, even when some transition entries vanish. The preparation modulus is part of the effectively compact data.

Box-cell integrals are computable uniformly: truncate an already certified tail and integrate the smooth bounded density on a compact box with its explicit derivative bounds. The complementary cell equals the normalization minus the interior cells. Simplex decoder nets are finite and rational. Matrix products and bounded linear audit costs have certified interval evaluations. This proves all of Definition~\ref{def:effective}. The general theorems apply. The proof never bounds a normalized likelihood ratio by a positive constant, and never reconstructs an exact posterior in the observer.
\end{proof}

The identity-transition member is a standard static hidden-state Gaussian sensing experiment, and positive-transition members are informative nonlinear filters. The theorem provides a two-sided finite certificate for their hard-state mean-budget optimum; it does not claim an explicit sharp exponent for all of them, nor an efficient solution of unknown-kernel learning. The parameter family is known to the program designer; the true parameter is not.

\subsection{Noncommuting quantum instruments, including singular effects}
Let $E$ be the real Hermitian matrices on $\mathbb C^q$, with the positive semidefinite cone, trace mass, and trace base norm. Let $Y$ be a compact metric report space with a fixed computable probability measure $\lambda$ and computable finite partitions of vanishing diameter. For each action and parameter, let $E_a^\theta(y)$ be positive semidefinite effects satisfying
\begin{equation}\label{eq:effective-quantum-effects}
 \int E_a^\theta(y)\lambda(dy)=I,\qquad
 0\le E_a^\theta(y)\le C_EI.
\end{equation}
Let $\mathcal T_a^\theta$ be a completely positive trace-preserving channel. The physical instrument is
\begin{equation}\label{eq:effective-quantum}
 L_a^\theta(y)\rho=
 \mathcal T_a^\theta\bigl(E_a^\theta(y)^{1/2}\rho E_a^\theta(y)^{1/2}\bigr).
\end{equation}
Take an effectively compact, uniformly computable family with computable uniform moduli for effects in $(\theta,y)$ and channels in the induced trace norm. Effects may lose rank. No simultaneous diagonalization is assumed. A protected informationally complete finite audit is fixed once and for all: a finite informationally complete POVM $\{F_b\}$ with real dual matrices $D_b$ satisfies $\sum_b\tr(F_b\rho)D_b=\rho$, and reports $D_b$. It has bounded, affine first and second moments. A finite such POVM can be obtained from sufficiently small positive perturbations of $I$ in a Hermitian basis and their negative partners; its finite spanning measurement matrix has a computable right inverse. Decoders lie in the density-operator body.

\begin{theorem}[Exact-budget quantum construction at changing rank]\label{thm:effective-quantum}
Under these raw hypotheses, including certified finite report partition execution, Theorem~\ref{thm:effective-budget} gives two-sided finite certificates for every rational $N\ge1$ and $M\ge2$. The observer is classical, its stopping and decoder labels are charged in $M$, and it uses no retained quantum memory. Rank-deficient preparations, rank-changing effects, and noncontracting unitary channels are allowed.
\end{theorem}
\begin{proof}
Complete positivity of \eqref{eq:effective-quantum} and integrated trace normalization follow from \eqref{eq:effective-quantum-effects}. The audit mean is $\rho$ and separates all density operators. Its finite set of first and second moments belongs to the Hermitian linear space; report pullbacks are linear in that same space. Thus the finite executable-test representation is verified without assuming a Fisher matrix or a smooth state parametrization.

For positive matrices $A,B$, the operator-norm square-root inequality
\begin{equation}\label{eq:sqrt-holder}
 \norm{A^{1/2}-B^{1/2}}_\infty\le\norm{A-B}_\infty^{1/2}
\end{equation}
is valid including at zero eigenvalues. One elementary proof puts $t=\norm{A-B}_\infty$ and uses $A\le B+tI$, operator monotonicity of the square root, and $(B+tI)^{1/2}\le B^{1/2}+\sqrt t I$; exchange $A,B$ for the opposite order. Operator monotonicity follows, for instance, by integrating $x/(x+s)$ against its positive square-root integral representation. These are finite-matrix statements.

For Hermitian $X$ with trace norm at most one, the ideal property of the trace norm and a two-term expansion give
\[
 \norm{A^{1/2}XA^{1/2}-B^{1/2}XB^{1/2}}_1
 \le(\norm{A^{1/2}}_\infty+\norm{B^{1/2}}_\infty)
                  \norm{A^{1/2}-B^{1/2}}_\infty.
\]
Trace-preserving positive maps contract trace norm on Hermitian inputs. Consequently effect and channel discrepancies $u,v$ give a raw instrument discrepancy at most
\begin{equation}\label{eq:quantum-raw-modulus}
 2\sqrt{C_E}\sqrt u+C_Ev.
\end{equation}
The same estimate in the report coordinate supplies a uniform oscillation bound on every partition cell, hence \eqref{eq:report-approximation}. No inverse of a density operator, normalized posterior, or effect is used. Continuity is therefore retained across rank changes.

Integrals of the computable uniformly continuous finite-dimensional matrix functions are computable with certified partition error under the stipulated computable measure. Density-operator decoder nets can be constructed by approximating a square-root factor with rational real and imaginary entries, forming $BB^*/\tr(BB^*)$, and taking sufficiently fine bounded factor grids; these are legal, finitely described positive matrices and cover the compact state body. The finite audit moments and all product evaluations are computable. Definition~\ref{def:effective} follows, so the general theorem proves the assertion.
\end{proof}

A concrete continuously indexed family takes $Y$ to be a compact symmetric trace-zero Hermitian ball, $E^\theta(H)=I+\eta_\theta H$ with $0\le\eta_\theta\le1$ and $\norm H_\infty\le1$, and any effectively compact channel family. When $\eta_\theta=1$ and $H$ has eigenvalue $-1$, an effect is singular. A symmetric computable measure on $Y$ has mean zero, so normalization holds. Discrete, Cantor-parametrized, or finite mixtures of such report measures also qualify when the specified partitions and integrals are computable. The actual weights remain in $\lambda$ and in every instrument integral; they are not replaced by ambient volume. Report support may vary with the incoming state through zero effects. A single dominating reference measure remains an assumption.

\subsection{The same normalized transfer in both realizations}
\begin{corollary}[Raw verification of the sharp-order transfer]\label{cor:raw-normalized}
Consider two Gaussian experiment families with the common compact parameter and report interfaces of Theorem~\ref{thm:effective-gaussian}. Suppose their preparations, transition probabilities, Gaussian means and covariance matrices differ by at most $u$, uniformly in the unknown model and action; retain the stated common ellipticity bounds. Then, for some constant $C$ depending only on these raw bounds, the same-budget risk difference is bounded by
\begin{equation}\label{eq:gaussian-normalized}
 \min\{L,\ L[2C(N+1)u+\sqrt{C^2(N+1)^2u^2+2c_NC(N+1)u}]+\alpha\}.
\end{equation}
Here $N\ge1,M\ge2$, and $\alpha$ denotes any additional uniform audit error.

For the quantum families of Theorem~\ref{thm:effective-quantum}, suppose preparations differ in trace norm by at most $u$, channels in the declared induced trace norm by at most $u$, and effects uniformly in operator norm by at most $u$. Put
\[
 z_u=u+CN(\sqrt u+u).
\]
The risk difference is at most
\begin{equation}\label{eq:quantum-normalized}
 \min\{L,\ L[2z_u+\sqrt{z_u^2+2c_Nz_u}]+\alpha\}.
\end{equation}
This conclusion permits rank changes. If both effect families instead have a common strictly positive lower bound, one may replace $\sqrt u+u$ by a constant multiple of $u$. These are sufficient moduli; the fractional effect-parameter exponent at rank change is not asserted sharp. Both realizations optimize over the entire declared common-program class at the same $N,M$.
\end{corollary}
\begin{proof}
For Gaussian densities, integrate the derivatives in the means and covariance entries along the line joining the two parameter values. Common ellipticity and bounded means bound these derivatives in $L^1$ by a constant. Multiplying by the transition probabilities and summing the finite input/output indices yields the integral induced base-norm instrument error $\eta\le Cu$. The preparation contributes at most $Cu$. The proof integrates over the whole report space; the tail is not discarded. Positivity implies the same bound after every common report-routing kernel. Apply Theorem~\ref{thm:normalized-transfer} with $\zeta\le C(N+1)u$.

In the quantum case, write $A=\sqrt E$, $B=\sqrt{\widetilde E}$. On a unit trace positive input $x$,
\[
 \|AxA-BxB\|_1\le(\|A\|+\|B\|)\|A-B\|.
\]
Use \eqref{eq:sqrt-holder}, the uniform effect bound and the probability reference measure, then add the channel error. Positive/negative decomposition extends the bound to the induced base norm, giving $\eta\le C(\sqrt u+u)$. This is an unnormalized-instrument estimate, not continuity of a posterior divided by a small likelihood. If the effects are bounded below by $\kappa I$, the identity
\[
 A(A-B)+(A-B)B=E-\widetilde E
\]
has the solution $A-B=\int_0^\infty e^{-tA}(E-\widetilde E)e^{-tB}\,dt$, whose operator norm is at most $u/(2\sqrt\kappa)$. This gives the improved linear bound. Theorem~\ref{thm:normalized-transfer} proves the assertions, including unchanged state and mean budgets and the separate audit error.
\end{proof}
